% Sources and the claim checked for each citation are recorded in
% literature-verification.md in this directory.
\subsection{Opinion dynamics and validation against survey data}\label{sec:bg-opinion}

Flache and colleagues distinguish assimilative influence, which reduces differences between connected agents; similarity-biased influence, which limits assimilation to agents who are already close; and repulsive influence, which can widen differences after contact \cite{flache2017}. Classic models average the opinions of contacts \cite{degroot1974}, keep an anchor to initial opinions \cite{friedkin1990}, or let agents compromise only with others within a confidence bound, in random pairs \cite{deffuant2000} or with all such neighbors at once \cite{hegselmann2002}; repulsive variants move agents away from sources that are too dissimilar \cite{jager2005}. Evidence for these shapes is mixed. Tak\'acs, Flache and M\"as found no support for the idea that discrepancy or disliking pushes opinions away from the source; shifts toward the source grew with opinion distance \cite{takacs2016}, which also speaks against a confidence bound. In field data on adolescents, a stochastic actor-oriented model found negative influence, but little support for bounded confidence, in preferences for one clothing style \cite{tang2025empirical}.

Many such models are judged by aggregate patterns, such as the segregated communities that a model of influence and unfollowing produced, consistent with Twitter data \cite{sasahara2021}. Answering calls for empirical tests of micro-level assumptions \cite{flache2017}, Gestefeld and Lorenz calibrate an interaction model to institutional-trust distributions from the European Social Survey and to annual trust transitions of SHP respondents \cite{gestefeld2023}. With repeated surveys and recorded communication among the same people, Nurek and colleagues find that a baseline in which everyone keeps their opinion reaches a high overall F1 score while detecting no change, and that on real-world data opinion-dynamics models often fail to outperform such trivial baselines \cite{coding2024}. Moor-Smith and Carpentras show with synthetic data that even the generating model, rerun with its true parameters, beats a no-change baseline only above a minimum amount of opinion change, and that where the signal is weak the calibrated model learns to mimic that baseline and ``freeze'' \cite{moor2026}. Random noise and scale distortions, unlike binning, strongly affect the predictions of the Deffuant model \cite{carpentras2022}, and monotone distortions of the scale can make data generated by that model look incompatible with it \cite{carpentras2023}. In the terms of pattern-oriented modeling \cite{grimm2005}, our design treats short-lag correlations as the pattern that competing processes must share and longer-lag correlations as the pattern that may tell them apart.

\subsection{Measurement error, stability and state dependence}\label{sec:bg-measurement}

Unstable answers have long been read in two ways. Converse's ``black-and-white'' model divided respondents into stable opinion holders and unstable changers whose answers are essentially random \cite{converse1964}; Achen attributed nearly all of the apparent inconsistency to vague questions and response categories \cite{achen1975}. Zaller and Feldman proposed that most people answer with whatever ideas are at the top of their heads \cite{zaller1992}, and averaging many items on the same issue removes much of the measurement error and reveals stable preferences \cite{ansolabehere2008}. In a four-wave Swiss panel, Hill and Kriesi \cite{hill2001,hillkriesi2001} added a group of durable changers to Converse's two groups, with results that undermine the measurement-error model. Reliability varies widely across items \cite{alwin2007,alwinkrosnick1991reliability}: of 293 core items of the General Social Survey, 133 had reliabilities above 0.70 and 15 fell below 0.40 \cite{hout2016}, and the average measurement quality of 67 European Social Survey questions ranges from 0.25 to 0.88 \cite{poses2021}.

Trait--state--error models \cite{kenny1995tse} add a stable trait to the classical panel models of stability and reliability \cite{heise1969,wiley1970}, but they often give improper or unstable estimates, especially when waves are few \cite{cole2005tso,ludtke2018starts}. Ignoring the trait has its own cost: under realistic assumptions the cross-lagged panel model is very likely to find spurious cross-lagged effects \cite{lucas2023,hamaker2015critique}, which is why trait-plus-autoregressive persistence (model B) is the main comparison here. In General Social Survey panels, most change looks like short-term change or measurement error \cite{kiley2020,vaiseykiley2021}, while in six long household panels, the SHP among them, persistent change over the life course accompanies substantial year-to-year fluctuation \cite{lersch2023}.

In economics and statistics the same problem appears as state dependence versus heterogeneity. Heckman distinguishes true state dependence, in which experience changes later choices, from spurious state dependence, in which past outcomes predict future ones only because they proxy for persistent unobserved differences \cite{heckman1981}. Earlier, Feller pointed out that the same distribution of event counts arises from contagion and from heterogeneity, so that deciding between them requires the correlation between time intervals, not the distribution alone \cite{feller1943}. Mover--stayer models mix people who never change category with people who move as a Markov chain \cite{blumen1955,goodman1961}. Persistence can also differ between people, as when the stability of party identification rises with age \cite{alwinkrosnick1991}, and sums of simple series, including a series measured with error, typically follow mixed autoregressive moving-average models \cite{granger1976}. Our persistence models cover autoregressive, trait, state-dependent and mover--stayer forms, but not several autoregressive speeds or speeds that vary between people.

\subsection{Identifying social influence}\label{sec:bg-identification}

Manski's identification problem does not vanish when reference groups are known: inference is then difficult or impossible if the variables that define groups and those that directly affect behavior are functionally dependent or statistically independent \cite{manski1993}. When interactions follow a known social network, endogenous and exogenous peer effects in an extended linear-in-means model are generally identified, although identification can fail for particular network structures \cite{bramoulle2009}.

Observational studies face a second confound: people form ties with similar others \cite{mcpherson2001}. Christakis and Fowler reported that obesity appears to spread through social ties \cite{christakis2007}; with the same kind of models, Cohen-Cole and Fletcher found ``network effects'' for acne, headaches and height that became smaller and insignificant after adjustment for environmental confounders \cite{cohencole2008}; matched-sample estimation attributed more than half of the apparent contagion in product adoption to homophily \cite{aral2009}, and sensitivity analysis asks how strong latent homophily or confounding would have to be to explain away a contagion effect \cite{vanderweele2011}. When ties and behavior are observed over several waves, stochastic actor-oriented models separate selection from influence \cite{snijders2010,steglich2010}, and simulations with known truth have tested estimators of influence and selection \cite{zachrison2016,daza2021}. Other work treats measurement error as error in the observed network rather than in the responses \cite{xu2021sensitivity,advani2018credibly}.

Some recent methods dispense with observed links, identifying linear network models from many separate networks with known membership \cite{lewbel2023} or from panels with many periods relative to the number of ties per unit \cite{depaula2024}, or inferring the network and the interaction kernel jointly from multiple trajectories \cite{lang2026joint}. None of these routes fits the panels studied here: each is a single realization with eleven waves, thousands of respondents and no known groups, and the influence rule is nonlinear.

Even with observed networks, influence is hard to detect: in simulated friendship networks with three-category opinions, a stochastic actor-oriented model detected an influence effect of 0.8 in 10\% of 200 data sets with two waves and 30 actors, against 97.5\% with five waves and 120 actors \cite{stadtfeld2020}.

\subsection{Simulation-based model discrimination}\label{sec:bg-discrimination}

The criterion used here is a classifier two-sample comparison. If two samples share a distribution, a classifier trained to tell them apart should stay near chance on held-out data; statistically significant accuracy above chance indicates that they differ in a way the classifier can use \cite{lopezpaz2017}. Such tests are consistent whenever the classifier's true accuracy stays bounded away from chance, and in a high-dimensional Gaussian setting a linear-discriminant test loses only a constant factor of efficiency relative to Hotelling's test \cite{kim2021classifier}. A classifier that sees a few summary statistics can still miss differences: in approximate Bayesian computation, model choice based on insufficient summaries can be unreliable \cite{robert2011}. Our persistence models are fitted by matching moments of simulated panels, in the spirit of indirect inference \cite{gourieroux1993}.

Two further results bound what such comparisons can show. Gelastopoulos, Sage and van de Rijt prove that the joint distribution of any stochastic process describing longitudinal observations can be reproduced exactly both by a pure reinforcement model, driven by past observations and independent noise, and by a pure heterogeneity model, driven by a fixed latent variable and noise \cite{gelastopoulos2025}. With functional forms left unrestricted, the theorem guarantees that a twin exists, not that a given parametric family can serve as one. In a preprint, Ure\~{n}a-Carrion, Peixoto and I\~{n}iguez show that likelihood-based selection among binary spreading mechanisms on known, fully observed networks can fail under common conditions \cite{urena2026}. A boundary against fixed persistence models matched at most at lags one and two is therefore a finite contrast that does not contradict generic equivalence.

\subsection{Positioning}\label{sec:bg-positioning}

Table~\ref{tab:positioning} compares the design with its closest antecedents. In the table, the studies that observe networks or communication do not model response error, the panel studies that model response error contain no influence between people, and the studies that fit influence models to survey data judge them by fit or against simple baselines, not against persistence matched to the same short-term stability. No single component is new; to our knowledge, based on a documented but limited search (\esmresource, Section~\ref{esm-sec:search}), the combination is.

% Two-column float with the structure of the class's table environment
% (8 bp body font, notes under the tabular), which table* lacks.
\begin{table*}[tp]
\begin{center}
\begin{threeparttable}
\tablebodyfont
\caption{Closest antecedents and the present design}\label{tab:positioning}
\setlength{\tabcolsep}{3pt}
\begin{tabular}{@{}p{0.17\textwidth}lllllp{0.21\textwidth}@{}}
\toprule
\raggedright Study & \begin{tabular}[t]{@{}l@{}}Network\\observed\end{tabular} & \begin{tabular}[t]{@{}l@{}}Exposure\\measured\end{tabular} & \begin{tabular}[t]{@{}l@{}}Response\\error\\modeled\end{tabular} & \begin{tabular}[t]{@{}l@{}}Individual\\trajectories\end{tabular} & \begin{tabular}[t]{@{}l@{}}Matched\\no-network\\persistence\end{tabular} & \raggedright Criterion\tabularnewline
\midrule
\raggedright Gestefeld and Lorenz \cite{gestefeld2023} & No & No & No & Partly$^{a}$ & No & \raggedright Fit to distributions and transition matrices\tabularnewline
\raggedright Nurek et al.\ \cite{coding2024} & Yes & Yes$^{b}$ & No & Yes & No$^{c}$ & \raggedright F1 and divergence against simple baselines\tabularnewline
\raggedright Moor-Smith and Carpentras \cite{moor2026} & No & No & No$^{d}$ & No$^{e}$ & No$^{c}$ & \raggedright Variance explained beyond no change\tabularnewline
\raggedright Kiley and Vaisey \cite{kiley2020} & No & No & Partly$^{f}$ & Yes & n/a$^{f}$ & \raggedright BIC of updating versus settled models\tabularnewline
\raggedright Mulder \cite{mulder2023}; Dang and Ma \cite{dang2026}, preprint & No & No & Yes & Yes & n/a$^{g}$ & \raggedright Power of a test of a within-person cross-lagged effect\tabularnewline
\raggedright Stadtfeld et al.\ \cite{stadtfeld2020} & Yes & Yes & No & Yes & No$^{h}$ & \raggedright Power of a test of the influence parameter\tabularnewline
\raggedright Ure\~{n}a-Carrion et al.\ \cite{urena2026}, preprint & Yes & Yes & No & Yes & No$^{i}$ & \raggedright Likelihood-based model selection\tabularnewline
\raggedright This study & No & No & Yes$^{j}$ & Yes & Yes$^{k}$ & \raggedright AUC of a fixed LDA on pooled lag-3--5 correlations\tabularnewline
\botrule
\end{tabular}
\par\vskip3pt\tablefootnotefont Network observed: ties among respondents are known. Exposure measured: what contacts say or do is recorded. Matched no-network persistence: a persistence process without a network, fitted to the same short-lag correlations. $^{a}$Transitions between consecutive annual waves. $^{b}$Communication events and survey opinions are observed; the opinion expressed in each event is simulated. $^{c}$Compared with simple baselines, such as unchanged opinions. $^{d}$Noise is added to the opinion process, not to responses. $^{e}$Repeated cross-sections. $^{f}$No influence model; settled dispositions and active updating are compared, and deviations from a settled disposition absorb measurement error together with short-lived change. $^{g}$No influence between people; the null is a zero within-person cross-lagged effect. $^{h}$Test of the influence parameter against zero in the same model. $^{i}$An independent spreading mechanism is among the candidate models. $^{j}$Known signal share $r$; eleven ordinal categories. $^{k}$Matched at lag one, or at lags one and two. BIC, Bayesian information criterion; AUC, area under the receiver operating characteristic curve; LDA, linear discriminant analysis.\par
\end{threeparttable}
\end{center}
\end{table*}
